\begin{table*}[t]\centering\caption{Main results (the 600 main runs, 20 seeds per cell): (a) exact-support rate, (b) mean and (c) median relative coefficient error, (d) mean coefficient error when the regression is restricted to the true support (isolates the regression problem from support selection).}\label{tab:main}
\small
\begin{tabular}{@{}c@{\hspace{1.5em}}c@{}}
(a) exact-support rate & (b) mean coefficient error\\[1pt]
\begin{tabular}{lrrrrrr}\toprule
System & 0\% & 0.5\% & 1\% & 2\% & 5\% & 10\% \\\midrule
FD & 1.00 & 1.00 & 0.95 & 0.80 & 0.00 & 0.00 \\
SG & 1.00 & 1.00 & 1.00 & 0.90 & 0.25 & 0.00 \\
Spline & 1.00 & 1.00 & 1.00 & 0.80 & 0.20 & 0.00 \\
TV & 1.00 & 1.00 & 1.00 & 0.85 & 0.40 & 0.00 \\
Weak & 1.00 & 1.00 & 1.00 & 1.00 & 0.65 & 0.10 \\
\bottomrule\end{tabular}
 & \begin{tabular}{lrrrrrr}\toprule
System & 0\% & 0.5\% & 1\% & 2\% & 5\% & 10\% \\\midrule
FD & 0.0065 & 0.0068 & 0.0118 & 0.0201 & 0.0469 & 0.3013 \\
SG & 0.0033 & 0.0037 & 0.0046 & 0.0126 & 0.0772 & 0.1721 \\
Spline & 1.5e-5 & 8.9e-4 & 0.0021 & 0.0150 & 0.0636 & 0.2485 \\
TV & 0.0025 & 0.0022 & 0.0020 & 0.0095 & 0.0536 & 0.1473 \\
Weak & 1.0e-8 & 8.6e-4 & 0.0017 & 0.0035 & 0.0337 & 0.0941 \\
\bottomrule\end{tabular}
\\[33pt]
(c) median coefficient error & (d) mean coefficient error on the true support\\[1pt]
\begin{tabular}{lrrrrrr}\toprule
System & 0\% & 0.5\% & 1\% & 2\% & 5\% & 10\% \\\midrule
FD & 0.0065 & 0.0071 & 0.0081 & 0.0129 & 0.0456 & 0.1337 \\
SG & 0.0034 & 0.0038 & 0.0046 & 0.0084 & 0.0739 & 0.1271 \\
Spline & 1.5e-5 & 7.7e-4 & 0.0019 & 0.0067 & 0.0571 & 0.1269 \\
TV & 0.0025 & 0.0021 & 0.0016 & 0.0046 & 0.0470 & 0.1224 \\
Weak & 1.1e-8 & 6.7e-4 & 0.0014 & 0.0027 & 0.0129 & 0.1094 \\
\bottomrule\end{tabular}
 & \begin{tabular}{lrrrrrr}\toprule
System & 0\% & 0.5\% & 1\% & 2\% & 5\% & 10\% \\\midrule
FD & 0.0065 & 0.0068 & 0.0078 & 0.0114 & 0.0351 & 0.1052 \\
SG & 0.0033 & 0.0037 & 0.0046 & 0.0082 & 0.0655 & 0.0988 \\
Spline & 1.5e-5 & 8.9e-4 & 0.0021 & 0.0058 & 0.0472 & 0.0987 \\
TV & 0.0025 & 0.0022 & 0.0020 & 0.0040 & 0.0324 & 0.0993 \\
Weak & 1.0e-8 & 8.6e-4 & 0.0017 & 0.0035 & 0.0094 & 0.0403 \\
\bottomrule\end{tabular}

\end{tabular}
\end{table*}

\begin{figure*}[t]\centering\includegraphics[width=0.72\textwidth]{noise_curves.pdf}
\caption{Exact-support rate (95\% Wilson intervals) and median coefficient error (interquartile range) versus noise level.}\label{fig:noise}\end{figure*}

\textit{H1 (weak form has a higher support rate than every baseline at every level from 1\%): refuted.} At 1\% the weak form, SG, spline and TV all recover every support (Table~\ref{tab:main}a), so there is no difference to test. At 2\% the weak form is at 1.00 against 0.80--0.90; the difference is significant against FD and spline ($p=\rhval{cmp/main/fd-stlsq/lorenz_n2/support_exact/welch_p:3}$, uncorrected) but not against SG or TV (Table~\ref{tab:pvals}a). At 5\% it is at 0.65 against 0.00--0.40, significant against FD, SG and spline but not against TV ($p=\rhval{cmp/main/tv-stlsq/lorenz_n5/support_exact/welch_p:3}$). At 10\% every system fails (weak 0.10, others 0.00, not significant). The weak form is never below a baseline, but the significant win over all four baselines that H1 requires holds at no level. The picture below 5\% depends on how the threshold is chosen (Sec.~\ref{sec:ablations}): with the first tie rule the weak form was at 0.75 at 0.5, 1 and 2\%, below baselines at 0.80--0.90.

\begin{table}[t]\centering\caption{Welch $p$-values ($^*$: $p<0.05$; --: identical rates). In (a) no baseline has a higher rate and in (b) the weak form has the lower mean in every cell; in (c) $^\ddagger$ marks a smoother whose mean is \emph{higher} than FD's.}\label{tab:pvals}
\begin{tabular}{lrrrrrr}\toprule
System & 0\% & 0.5\% & 1\% & 2\% & 5\% & 10\% \\
\midrule\multicolumn{7}{l}{(a) support rate, weak form vs.} \\
FD & -- & -- & 0.330 & 0.042$^*$ & 0.000$^*$ & 0.163 \\
SG & -- & -- & -- & 0.163 & 0.010$^*$ & 0.163 \\
Spline & -- & -- & -- & 0.042$^*$ & 0.003$^*$ & 0.163 \\
TV & -- & -- & -- & 0.083 & 0.119 & 0.163 \\
\midrule\multicolumn{7}{l}{(b) coefficient error, weak form vs.} \\
FD & 0.000$^*$ & 0.000$^*$ & 0.024$^*$ & 0.008$^*$ & 0.179 & 0.037$^*$ \\
SG & 0.000$^*$ & 0.000$^*$ & 0.000$^*$ & 0.040$^*$ & 0.000$^*$ & 0.004$^*$ \\
Spline & 0.000$^*$ & 0.855 & 0.287 & 0.055 & 0.007$^*$ & 0.012$^*$ \\
TV & 0.000$^*$ & 0.000$^*$ & 0.550 & 0.170 & 0.088 & 0.014$^*$ \\
\midrule\multicolumn{7}{l}{(c) coefficient error, FD vs.} \\
SG & 0.000$^*$ & 0.000$^*$ & 0.096 & 0.286 & 0.000$^*$$^\ddagger$ & 0.187 \\
Spline & 0.000$^*$ & 0.000$^*$ & 0.030$^*$ & 0.526 & 0.004$^*$$^\ddagger$ & 0.626 \\
TV & 0.000$^*$ & 0.000$^*$ & 0.027$^*$ & 0.137 & 0.369$^\ddagger$ & 0.116 \\
\bottomrule\end{tabular}
\end{table}

\textit{Where the failures come from.} Table~\ref{tab:fpfn} separates spurious terms (FP) from missed terms (FN). Up to 2\% the weak form has neither. The baselines at 2\% both add (0.10--0.20 per trial) and miss (0.05--0.10) terms. At 5\% the baselines mostly add terms (0.65--2.70 FP against 0.00--0.25 FN), while the weak form has 0.20 FP and 0.25 FN. FD at 5\% never recovers the support and carries 2.70 spurious terms per trial at its threshold of 0.1, picked by median error alone (Sec.~\ref{sec:setup}); larger thresholds do not give a high rate either (0.15 at 0.4 and 0.25 at 0.6, Table~\ref{tab:thr}).

\begin{table}[t]\centering\caption{Mean false positives / false negatives per trial (30 coefficients, 7 nonzero).}\label{tab:fpfn}
\setlength{\tabcolsep}{4pt}\begin{tabular}{lrrrrr}\toprule
Noise & FD & SG & Spline & TV & Weak \\\midrule
0\% & 0.00/0.00 & 0.00/0.00 & 0.00/0.00 & 0.00/0.00 & 0.00/0.00 \\
0.5\% & 0.00/0.00 & 0.00/0.00 & 0.00/0.00 & 0.00/0.00 & 0.00/0.00 \\
1\% & 0.00/0.05 & 0.00/0.00 & 0.00/0.00 & 0.00/0.00 & 0.00/0.00 \\
2\% & 0.20/0.10 & 0.10/0.05 & 0.15/0.10 & 0.15/0.05 & 0.00/0.00 \\
5\% & 2.70/0.00 & 0.85/0.15 & 0.85/0.20 & 0.65/0.25 & 0.20/0.25 \\
10\% & 1.45/1.20 & 2.85/0.80 & 1.30/1.25 & 2.60/0.70 & 0.75/0.85 \\
\bottomrule\end{tabular}
\end{table}

\textit{H2 (each smoother has a lower mean coefficient error than FD from 1\%): refuted.} At 1, 2 and 10\% the means of SG, spline and TV are below FD (Table~\ref{tab:main}b), significantly only for spline and TV at 1\% (Table~\ref{tab:pvals}c). At 5\% all three are above FD (0.0536--0.0772 against 0.0469), significantly for SG and spline, which refutes H2 as registered. This verdict depends on the FD threshold: at the first-version value of 0.05 the FD error is 0.1354, above all three smoothers at their tuned thresholds (Table~\ref{tab:threrr}), and 0.05 and 0.1 were tied in the tuner although they give different models on three of the five tuning trajectories (Sec.~\ref{sec:setup}). The claim for TV is the weakest: its higher mean is not significant (Table~\ref{tab:pvals}c), and at equal thresholds TV is below FD in all five columns of Table~\ref{tab:threrr}. Spline is above FD at four of the five equal thresholds (not at 0.05) and SG at all five; on the true support FD is below SG and spline but not TV (0.0351 against 0.0655, 0.0472 and 0.0324, Table~\ref{tab:main}d).

\textit{H3 (all systems exact at 0\% noise): supported.} All five systems recover the support in all 20 trials at 0\% and at 0.5\%. Under the first tie rule ($\lambda=0.05$) FD and SG were at 0.95, a threshold effect: both are at 1.00 for every $\lambda\ge0.1$ (Table~\ref{tab:thr}).

\textit{Coefficient error.} The weak form has the lowest mean coefficient error at every noise level (Table~\ref{tab:main}b) and the lowest error on the true support at every level (Table~\ref{tab:main}d; 0.0094 against 0.0324--0.0655 at 5\%, 0.0403 against 0.0987--0.1052 at 10\%). The differences in Table~\ref{tab:main}b are significant against all four baselines only at 0\% and 10\% (Table~\ref{tab:pvals}b); against spline they are not significant from 0.5 to 2\%, against TV not from 1 to 5\%, and against FD not at 5\%. At 10\% the gap is much smaller in the medians (0.1094 against 0.1224--0.1337, Table~\ref{tab:main}c) than in the means (0.0941 against 0.1473--0.3013). The clearest separation is at 5\%, where the weak-form median is 0.0129 against 0.0456--0.0739 (Fig.~\ref{fig:noise}).

\textit{What the results do not show.} The ordering of the systems below 5\% noise is not resolved: it reverses with the tie rule of the tuner. Rate differences are poorly resolved at $n=20$: \rhval{cmp/main/fd-stlsq/lorenz_n2/support_exact/delta:2} (1.00 against 0.80 at 2\%) gives $p=\rhval{cmp/main/fd-stlsq/lorenz_n2/support_exact/welch_p:3}$, while \rhval{cmp/main/tv-stlsq/lorenz_n2/support_exact/delta:2} at the same level ($p=\rhval{cmp/main/tv-stlsq/lorenz_n2/support_exact/welch_p:3}$) and \rhval{cmp/main/tv-stlsq/lorenz_n5/support_exact/delta:2} at 5\% (0.65 against 0.40, $p=\rhval{cmp/main/tv-stlsq/lorenz_n5/support_exact/welch_p:3}$) are not significant. The $p$-values are uncorrected over the 18 to 24 comparisons of each table; the two $p=\rhval{cmp/main/fd-stlsq/lorenz_n2/support_exact/welch_p:3}$ cells would not survive a Bonferroni correction even over the four comparisons of that level.
